\subsection{ABSOLUTELY SUMMABLE FAMILIES IN A NORMED SPACE}

DEFINITION 8. In a normed space E, a family \((\pmb{x}_{t})\) of points of E is said to be absolutely summable if the family \((||\pmb{x}_{t}||)\) of norms of the \(\pmb{x}_{t}\) is summable in R.

This concept appears to depend on the norm chosen on E; but by Proposition 7 of no. 3 and the comparison principle for summable families of real numbers, a family which is absolutely summable with respect to a norm p on E is absolutely summable with respect to any norm on E which is equivalent to p.

If \((\pmb{x}_t)_{i\in I}\) is a family of points of \(E\) which is summable and absolutely summable, we have

\[
\left| \left| \sum_ {i \in I} x _ {i} \right| \right| \leqslant \sum_ {i \in I} | | x _ {i} | |.\tag{4}
\]

Indeed, for each finite subset J of I we have \(\left\|\sum_{i\in J}x_{i}\right\| \leqslant \sum_{i\in J}||x_{i}||\) , and the inequality (4) follows by passing to the limit with respect to the directed set of finite subsets of I.

PROPOSITION 11. In a complete normed space E, every absolutely summable family is summable.

For if \((\pmb{x}_i)\) is an absolutely summable family in \(\mathbf{E}\) , then for each \(\varepsilon > 0\) there is a finite subset \(J\) of the index set \(I\) such that, for each finite subset \(H\) of \(I\) which does not meet \(J\) , we have \(\sum_{i \in H} ||\pmb{x}_i|| \leqslant \varepsilon\) ; hence a fortiori \(\left\| \sum_{i \in H} \pmb{x}_i \right\| \leqslant \varepsilon\) , and this proves the proposition, since \(E\) is complete (Cauchy's criterion, Chapter III, § 5, no. 2, Theorem 1).

A series whose general term is \(x_{n}\) is said to be absolutely convergent in E if the series whose general term is \(\|x_{n}\|\) is convergent in R, or (equivalently) if the family \((x_{n})\) is absolutely summable; consequently (Chapter III, § 5, no. 7, Proposition 9):

COROLLARY. In a complete normed space E, every absolutely convergent series is commutatively convergent.

The converse of Proposition 11 is in general false.

Consider for example the space \(\mathcal{B}(\mathbf{N};\mathbf{R})\) of bounded sequences \(x = (x_{n})_{n\in \mathbb{N}}\) of real numbers, with the norm \(\| x\| = \sup_n |x_n|\) . Let \(x_{m}\) be the sequence \((x_{mn})_n\in \mathbb{N}\) such that \(x_{mn} = 0\) if \(m\neq n\) and \(x_{mm} = 1 / m\) for \(m\geqslant 1\) . It is immediately verified that the sequence \((x_m)_{m\in \mathbb{N}}\) is summable in \(\mathcal{B}(\mathbf{N};\mathbf{R})\) and that its sum is the element \(y = (y_{n})\) such that \(y_0 = 0\) and \(y_{n} = 1 / n\) if \(n\geqslant 1\) ; but since \(\| x_m\| = 1 / m\) , the sequence of norms of the \(x_{m}\) is not summable in \(\mathbf{R}\) .

However, we have seen in Chapter VII, § 3, no. 1, that every summable family in \(\mathbf{R}^n\) is absolutely summable.

